{\small\textbf{March 2020 (Zheng, 1:30) [1:00–2:30]}}\par\smallskip
Let's start in March 2020. COVID has just hit Europe.
\begin{itemize}\setlength\itemsep{0pt}
\item Within days, Denmark, Norway, Sweden, then Germany, France, Ireland, Lithuania and Belgium \textbf{release} their countercyclical buffers.
\item \textcolor{navy}{\textbf{[def]}} A capital buffer is extra equity that banks must hold above the legal minimum, so they can absorb losses without failing.
\item Releasing it tells banks: use this capital now, keep lending.
\end{itemize}
\par\smallskip
The Netherlands could not do that. Its countercyclical buffer was at \textbf{zero}, so there was \textbf{nothing to release}.
\par\smallskip
So DNB improvised.
\begin{itemize}\setlength\itemsep{0pt}
\item On 17 March 2020 it cut the systemic buffers of its three largest banks, ING, Rabobank and ABN AMRO, from three percent to 2.5, 2 and 1.5 percent.
\item It also postponed a planned floor on mortgage risk weights.
\item Together that freed about \textbf{eight billion euros} of capital, which DNB said could support up to \textbf{200 billion} of lending.
\end{itemize}
\par\smallskip
And then DNB made a promise \emph{(pause)}: once things were back to normal, it would rebuild this capital as a \textbf{two-percent countercyclical buffer}. Keep that promise in mind; the whole story turns on it.
